\part{\trArnoldConjecture{}与\trFloerHomology}%

\chapter*{第二部分引言}
\addcontentsline{toc}{chapter}{第二部分引言}
\markboth{第二部分引言}{第二部分引言}%

本书第二部分致力于\textbf{\trFloerHomology}（Floer homology）的构造，其目标在于证明\textbf{\trArnoldConjecture}（Arnold conjecture）。这条猜想断言：定义在\trSymplectic\trManifold{} $W$ 上的一个\trHamiltonian\trVectorField，其\trPeriodic{}为 $1$ 的\trPeriodic\trTrajectory{}数目不小于\footnote{译注：原书此式即照此排印。按\trMorse{}不等式的口径，和式各项应取其维数，即 $\sum_k\dim HM_k(W;\ZZ/2)$；而且这一形式的\trArnoldConjecture{}（其精确陈述见第 6 章）还须假定 $W$ \trCompact{}无边界、所用\trHamiltonian{}量\trNondegenerate，原书此处均省而未言。——译者注}
\[
\sum_k HM_k(W;\ZZ/2)\text{。}
\]

正如我们在本书前言中所说，\trMorse{}不等式与\trArnoldConjecture{}的这种并置绝非偶然。事实上，这一证明策略（归功于\trFloer）在于把所考虑的那些\trTrajectory{}呈现为\trCriticalPoint——不过不再是一个定义在\trManifold{}上的函数的\trCriticalPoint，而是一个定义在 $W$ 的\trLoopSpace{}上的“\textbf{\trActionFunctional}（action functional）”\footnote{自\trHadamard{}以来，如此命名定义在无穷维空间上的函数已成传统。}的\trCriticalPoint（一条\trPeriodic\trTrajectory{}其实就是 $W$ 的一条\trLoop）。下一步是构造一个与\trMorse{}的那个类似的\trComplex——即\textbf{\trFloerComplex}（Floer complex）——其\trDifferential{}“计数”\textbf{\trFloerEquation}（Floer equation）的解。由构造本身可知，这个\trComplex{}的\trHomology——人们称之为\trFloerHomology——与\trPeriodic\trTrajectory{}的数目相关联。我们接下来还需证明：这一\trHomology{}与该\trManifold{}的\trMorseHomology\trIsomorphic。

我们将在\trArnoldConjecture{}证明纲要的展开处给出这一策略的细节，并在本书这一部分中加以应用。让我们在此既给出\trArnoldConjecture{}证明的一个概要，也给出这一部分内容的一个总览。

我们要做的第一件事，就是构造“\trActionFunctional”（在\cref{sec:action-functional}中）——它定义在我们\trManifold{}的一个\trLoopSpace{}上，其\trCriticalPoint{}正是我们所要寻找的那些\trPeriodic{}解。这就是\trMorseFunction{}的类似物。有了它的诸\trCriticalPoint，我们将尝试像在\trMorse{}的情形中那样构造一个\trComplex。由诸\trCriticalPoint{}所生成的、取在 $\ZZ/2$ 上的那个\trVectorSpace{}，其\trGraded{}由\textbf{\trMaslovIndex}（Maslov index）给出——它将在第 7 章中得到定义。

为了定义\trDifferential，我们将使用\trPseudogradient{}场的一个类似物——此处即\trActionFunctional{}相对于某个\trMetric{}的（负）\trGradient（参见\cref{sec:gradient-floer-equation}）。我们将证明：它那些\trEnergy{}有限的\trTrajectory——“\trFloerEquation”的解——连接着两个\trCriticalPoint。与\trMorse{}的情形一样，我们将对诸\trTrajectory{}的空间证明一个\trCompact{}性的结果（\cref{thm:gromov-compactness-theorem}）。一般说来，凭借\trTransversality{}方面的性质，这些\trTrajectory{}的空间都是\trManifold（第 8 章）。这就使得我们得以把这些\trTrajectory{}当作一个零维\trManifold{}中的元素来计数，并从而定义\trFloerComplex{}的\trDifferential。

为了证明后者确实是一个\trComplex——换言之，$\partial\circ\partial=0$——我们将像在\trMorse{}的情形中那样建立一个\trGluing{}性质（第 9 章）。同样与\trMorse{}的情形一样，我们将在第 11 章中证明：所得\trComplex{}的\trHomology{}既不依赖于泛函的选取，也不依赖于\trVectorField{}的选取。在此之前，我们将在第 10 章中证明：对一个不依赖于时间的 $\cC^2$-小\trHamiltonian{}量，\trFloerHomology{}与该\trManifold{}的\trMorseHomology{}相重合。因此，这在一般情形下也将成立。于是，这个\trComplex{}中诸\trVectorSpace{}的维数与其\trHomology{}的维数之间那个显而易见的不等式，便会给出所期待的结果，其方式与\trMorse{}不等式完全类似（\cref{thm:morse-inequalities}）。

于是，那个曾充当本书第一部分之研究对象的\trMorseComplex，便为第二部分充当起向导。如今所研究的对象要复杂得多，需要一种精微的分析。\trFloerEquation——一个关于\trActionFunctional{}的“\trGradient\trTrajectory”的\trDifferential{}方程——是一个偏\trDifferential{}方程，我们得以把\trEllipticRegularity{}技巧应用于它（将在第 12 章中明确给出）。这个\trComplex{}的定义需要这样一个事实：该方程的\trLinearization{}给出\trFredholmOperator。而那些使我们得以断言\trFloerComplex{}确实是一个\trComplex{}的性质，其证明在技术上远比\trMorseComplex{}的相应陈述之证明繁难。

我们把一批技术性的引理归拢到了单独的一章之中（\LaTeX{}已明智地把这一章编为第 13 章，迷信的读者尽可跳过），这些引理尤其是\trGluing{}性质的证明所要用到的。
